\section{What should not be automated, and the risks}

Three parts of the workflow stay with people in the proposed system: the question and the storyline of a scenario, the acceptance of the calibration, and the interpretation of the results.
This is not a precaution to be relaxed once the agents prove themselves.
Each is a judgment for which there is no correct answer to check an agent against: a storyline is a view of how the market could develop, a calibration is one of several fits that reproduce the reference values (Section 2.6), and an interpretation decides what a study shows.
Calibrating economic models to a benchmark is an established method precisely because the modeller takes responsibility for these choices \citep{kydland1996computational}.
An agent can prepare each of these decisions, but it cannot be accountable for them.

Keeping a person at these points only helps if the person actually judges.
People tend to rely on automated advice more than they should, and to check it less the more often it has been right \citep{parasuraman2010complacency}.
The system is designed to be right most of the time, so its success invites exactly this complacency.
An analyst who has accepted a hundred well-reasoned suggestions will read the hundred-and-first less carefully.
Two features of the design work against this.
The confidence classes direct attention to the values that most need it, so that careful review is asked for where it matters rather than everywhere (Section 5.6).
And the reasons attached to every suggestion give the analyst something to disagree with, not only a number to accept.
Neither removes the risk.
It also means that the measure proposed in Section 5.9, the share of suggestions accepted unchanged, is ambiguous: a rising share may mean that the agents have learned, or that the analyst has stopped looking.
The master's thesis should therefore also track how often analysts edit or reply, and check a sample of accepted suggestions independently.

A second risk is the loss of skill.
Automating the routine parts of a job leaves people with the hardest parts and less practice at them \citep{bainbridge1983ironies}, and the same has been observed for generative AI \citep{simkute2024ironies}.
If the calibration agent does the diagnosis, the analyst gets less practice at diagnosing, which is the skill needed to judge whether the agent's diagnosis is right.
For GGM this matters more than usual, because the knowledge of how to calibrate the model is held by few people and is largely unwritten (Sections 3.5 and 5.6).
New analysts would gain most from the system, as in other settings where AI helps the least experienced most \citep{brynjolfsson2025generative}, but they are also the ones least able to notice when it is wrong.
The lessons partly answer this, because they write down rules that today exist only in the developers' heads, and an analyst can read them.
Calibrating a scenario by hand from time to time, and comparing the result with the agent's, would keep the skill alive and test the agent at the same time.

A third risk is the task that looks routine but is not.
AI improves work on tasks within its capabilities and makes it worse on tasks just outside them that look similar \citep{dellacqua2026navigating}.
Looking up an official gas price looks like any other lookup until the country subsidises domestic gas.
The lessons cover the cases that have been found; the risk lies in those that have not, and the analyst's review is the only protection against them.
